
\thm{Kunita-Watanabe}{
    Let $M \in \mathcal{M}_0^2$ and $H \in L^2(M)$. Then $\forall~N \in \mathcal{M}_0^2$,
    $$
    \mathbb{E}\left[ 
        (H \cdot M)_{\infty} \cdot N_\infty
     \right] = \mathbb{E}\left[ 
        (H \cdot [M,N])_{\infty}
      \right].
    $$
    Moreover, 
    $$
        [H\cdot M , N] = H \cdot [M,N] \xlongequal{def}\int H_s \diff [M,N]_s.
    $$
}

\thmnn{
    Let $M,N \in \mathcal{M}_0^2$. Then 
    $$
    \left( \int_0^\infty \abs{ H_s K_s } \diff [M,N]_s \right)^2 \leq \int_0^\infty H_s^2 \diff [M]_s \cdot \int_0^\infty K_s ^2 \diff [N]_s. \quad \text{a.s.}
    $$
}

\pf{
    $\abs{ H_s K_s } = H_s K_s \cdot \mathrm{sgn}(H_s K_s)$.

    $\abs{ \diff [M,N]_s } = \diff [M,N]_s \cdot \mathrm{sgn}([M,N]_s)$.
    $$
    \begin{aligned}
        LHS &= \int_0^\infty H_s K_s \cdot \mathrm{sgn}(H_s K_s) \cdot \mathrm{sgn} ([M,N]_s) \diff [M,N]_s\\
        &=\int_{0}^{\infty} \tilde{H_s}K_s \diff [M,N]_s
    \end{aligned}
    $$
    $$
    \begin{aligned}
        RHS &= \int_0^\infty \tilde{H_s} \diff [M]_s \cdot \int_0^\infty K_s^2 \diff [N]_s
    \end{aligned}
    $$
    $$
    \Leftrightarrow \left( \int_0^\infty \tilde{H_s} K_s \diff [M,N]_s  \right)^2 \leq \int_0^\infty \tilde{H_s} \diff [M]_s \int_0^\infty K_s ^2 \diff [N]_s.
    $$
    $$
    \Leftrightarrow \begin{aligned}
        & H_n = \sum_{i=1}^n U_i (t_i,t_{i+1}]\\
        & K_n \equiv \sum_{i=1}^n V_i(t_i,t_{i+1}].
    \end{aligned}
    $$
    $$
    \Leftrightarrow \left( \sum_{i=1}^{n} U_i V_i \left( [M,N]_{t_{i+1}} - [M,N]_{t_i} \right) \right)^2 \leq \sum_{i=1}^2 U_i^2([M]_{t_{i+1}} - [M]_{t_i} ) \cdot \sum_{i=1}^n V_i^2 ([N]_{t_{i+1}}-[N]_{t_i}).
    $$

    It suffices to prove
    $$
    \abs{ [M,N]_{t_{i+1}} - [M,N]_{t_i} } \leq \left( [M]_{t_{i+1}} - [M]_{t_i} \right)^{\frac{1}{2}} \left( 
        [N]_{t_{i+1}} - [N]_{t_i}
     \right)^{\frac{1}{2}}.
    $$
    like $(\sum_i a_i b_i)^2 \leq \sum_i a_i^2 \sum_i b_i^2.$

    To prove this, note that $\forall~ \lambda \in \mathbb{Q}$, $[M+ \lambda N]$ is increasing and so 
    $$
        [M+ \lambda N]_{t_{i+1}} - [M+\lambda N]_{t_i} \geq 0 \quad \text{a.s..}
    $$
    
    Also 
    $$
        [M + \lambda N]_t = [M]_t + 2\lambda[M,N]_t + \lambda^2 [N]_t.
    $$
    because
    $
    (M+ \lambda N)_{t}^2 - [M+ \lambda N]_t
    $ is U.I. and by the uniqueness (?)
    $$
    M_t^2  + 2 \lambda M_t N_t + \lambda^2 N_t^2 -([M]_t + 2 \lambda [M,N]_t + \lambda^2 [N]_t).
    $$
    Hence 
    $$
    \left( [M]_{t_{i+1}} - [M]_{t_i} \right) + 2 \lambda \left( [M,N]_{t_{i+1}} - [M,N]_{t_i} \right) + \lambda^2 \left( [N]_{t_{i+1}} -[N]_{t_i} \right) \geq 0
    $$
    $$ 
    \Leftrightarrow b^2 - 4ac \leq 0.
    $$
    So
    $$
        \left( [M,N]_{t_{i+1}} - [M,N]_{t_i} \right)^2 \leq \left( [M]_{t_{i+1}} - [M]_{t_i} \right)\left( [N]_{t_{i+1}} - [N]_{t_i} \right).
    $$
    $$
        \mathbb{E}\left[ 
            (H\cdot M)_{\infty} N_{\infty}
         \right] = \mathbb{E}\left[ 
            (H \cdot [M,N])_{\infty}
          \right].
    $$

    If $H = Z(S,T]$, 
    $$\begin{cases}
     LHS = \mathbb{E}\left[ Z(M_T - M_S) \cdot N_\infty \right]\\
     RHS = \mathbb{E}\left[ Z([M,N]_{T} - [M,N]_S) \right]   
    \end{cases}$$

    $$
    \begin{aligned}
        LHS &= \mathbb{E}\left[ \mathbb{E}(N_\infty \mid \mathcal{F}_T) \cdot Z(M_T - M_S) \right] \\
        &= \mathbb{E}\left[ Z(M_T N_T - M_S N_T) \right]\\
        &= \mathbb{E}[ZM_TN_T] - \mathbb{E}\left[ \mathbb{E}(N_T \mid \mathcal{F}_S)\cdot ZM_S \right]\\
        &=\mathbb{E}\left[ Z M_T N_T \right] - \mathbb{E}[ZM_S N_S].
    \end{aligned}
    $$
    
    Because $M_t N_t - [M,N]_t$ is U.I.
    $$
    \mathbb{E}[M_T N_T] = \mathbb{E}[[M,N]_T] \implies \mathbb{E}[Z(M_T N_T - [M,N]_T)] = 0.
    $$
    And
    $$
    \begin{aligned}
        LHS &= \mathbb{E}\left[ Z(M_T N_T - M_S N_S) \right] \\
        &= \mathbb{E}\left[ Z \mathbb{E}\left[ M_TN_T -M_S N_S \mid \mathcal{F}_S \right] \right]\\
        &= \mathbb{E}\left[ Z \mathbb{E}[[M,N]_T - [M,N]_S \mid \mathcal{F}_S] \right]\\
        &=\mathbb{E}[Z([M,N]_T - [M,N]_S)].
    \end{aligned}
    $$
    Hence $\forall~ H \in b \varepsilon$
    $$
    \mathbb{E}[(H\cdot M)_{\infty} \cdot N_{\infty}] = \mathbb{E}\left[ (H \cdot [M,N])_{\infty} \right].
    $$
    
    If $H_n \xlongrightarrow{L^2(M)}H$, i.e.
    $$
    \mathbb{E}\left( \int_0^\infty (H_n(s) - H_s)^2 \diff [M]_s \right) \to 0
    $$
    $\implies (H_n \cdot M)_{\infty} \xlongrightarrow{L^2} (H\cdot M)_{\infty}$ by 
    $$
    \mathbb{E}[(H_n \cdot M - H\cdot M)_{\infty}^2] = \mathbb{E}\left[ \int_{0}^{\infty}(H_n - H) \diff [M]_s \right] \to 0.
    $$
    
    Hence $\mathbb{E}[(H_n \cdot M)_{\infty} \cdot N_{\infty}] \to \mathbb{E}[(H\cdot M)_{\infty} N_\infty]$.

    Next,
    $$
    \begin{aligned}
    \mathbb{E}\left[ \abs{ (H_n \cdot [M,N])_{\infty} - (H \cdot [M,N])_{\infty} } \right] &= \mathbb{E}\abs{ ((H_n - H) \cdot [M,N])_{\infty} } \\
    &\leq \left( \mathbb{E}\left( \int_0^\infty(H_n -H) \diff [M]_s \right) \right)^{\frac{1}{2}}\left( \mathbb{E}[N]_{\infty} \right)^{\frac{1}{2}} \to 0
    \end{aligned}
    $$
    by using 
    $$
    \left( \int_0^\infty \abs{ H_s } \diff [M,N]_s \right) \leq \left( \int_0^\infty H_s^2 \diff [M]_s \right)^{\frac{1}{2}} \left( \int_0^\infty 1 \diff [N]_s \right)^{\frac{1}{2}}.
    $$

    Hence $\mathbb{E}[(H_n \cdot [M,N])_{\infty}] \to \mathbb{E}[(H\cdot [M,N])_{\infty}]$.
}

\rmkb{
    proof of $[H\cdot M,N] = H \cdot [M,N]$
    $\Leftrightarrow$ It suffices to. prove $(H\cdot M)_{t} N_t - \int_{0}^\infty H_s \diff [M,N]_s$ is a U.I. martingale.
    
    $\Leftrightarrow \forall~$stopping time $R$
    $$
    \mathbb{E}[(H\cdot M)_{R} \cdot N_R] = \mathbb{E}\left[ \int_0^R H_s \diff [M,N]_s \right].
    $$
    $$
    \begin{aligned}
        LHS &= \mathbb{E}\left[ (H\cdot M)_{\infty}^R \cdot N_\infty^R \right] \\
        &= \mathbb{E}\left[ 
        (H^R \cdot M)_{\infty} \cdot N_\infty^R
        \right] \\
        &= \mathbb{E}[(H^R \cdot [M,N^R])_{\infty}]\\
        &= \mathbb{E}\left[ \int_0^\infty H^R(s) \diff [M,N^R]_s \right]\\
        &= \mathbb{E}\left[ \int_0^R H(s) \diff [M,N]_s \right] = RHS.
    \end{aligned}
    $$
}

\defn{local martingale}{
    $\mathcal{M}_{loc}=\{\text{local martingale}\}$. 
    
    $\exists~$stopping times $T_n \uparrow \infty$ s.t. $M^{T_n}$ is a martingale.
}
\defn{continuous local martingale}{
    $c\mathcal{M}_{loc}=\{\text{continuous local martingale}\}$.

    Quadratic variation of $c \mathcal{M}_{loc}$. 
    
    If $M \in c\mathcal{M}_{loc}$, then $\exists~[M]$ increasing process such that $M^2 - [M] \in c \mathcal{M}_{loc}$.
}

\lemnn{
    If $M \in c \mathcal{M}_{loc}$. Fix $n \geq 1$, define
    $$
    T_n = \inf \{t \geq 0\colon \abs{ M_t }\geq n\}.
    $$
    Then $M^{T_n}$ is a bounded martingale.
}

\pf{
    By definition, $\exists~$stopping $(S_k)\uparrow \infty$ s.t. $M^{S_k}$ is a martingale.

    $\implies M^{S_k \wedge T_n}$ is also a martingale.

    So $\forall~ s < t$, $$\mathbb{E}(M^{S_k \wedge T_n}_t \mid \mathcal{F}_s) = M_s^{S_k \wedge T_n}.$$

    Let $k \to \infty$ to get $\mathbb{E}(M_t^{T_n} \mid \mathcal{F}_s) = M_s^{T_n}.$ 

    pf of Quadratic variation:

    \textbf{(Existence.)} $\forall~ n \geq 1$, we let
    $$
    T_n = \inf \{t \geq 0 \colon \abs{ M_t } \geq n\}.
    $$
    Then define $\forall~ t < T_n$,
    $$
    [M]_t = [M^{T_n}]_t .
    $$

    We need to check that $[M^{T_n}]_t$ is well-defined.
    $$
    \begin{aligned}
        [M^{T_2}]^{T_1} = [M^{T_2}]_{t\wedge T_1}
    \end{aligned}
    $$
    Because $[M]^T = [M^T]$,
    $$
    [M^{T_2}]^{T_1} = [M^{T_2 \wedge T_1}] =[M^{T_1}].
    $$
    By $T_n \uparrow \infty$, we have $[M]_t$ is defined for all $t \geq 0$.

    \textbf{(uniqueness.)} If $\exists~A,A'$ increasing such that $N = M^2 - A$ and $N' =M^2 - A' \in c \mathcal{M}_{loc}$.
    Then $N-N'= A' - A \in c \mathcal{M}_{loc}$ is of finite variation.
    
    \lemnn{
        If $M \in c \mathcal{M}_{loc}$ and $M$ is of finite variation, then $M_t \equiv 0.$
    }
    \pf{
        Let $V_M(t) =$ total variation of $M$ on $[0,t]$. Define $T_n = \inf \{t \geq 0 \colon V_M(t) > n\}.$ 

        Then $M^{T_n}$ is a bounded martingale and is of finite variation. Next, define 
        $$
        A_t^n \equiv \sum_{k=1}^{2^n}\left( M(t_{k}^n) - M(t_{k-1}^n) \right)^2
        $$
        with $t_k^n = \frac{k}{2^n} t$.

        So
        $$
        A_t^n \leq \sup_{1 \leq k \leq 2^n} \abs{ M(t_k^n) - M(t_{k-1}^n) }\cdot V_M(t) \xlongrightarrow{a.s.} 0
        $$
        by continuity of $M$.

        Thus $A_t^n \to 0$, $\forall~ t \geq 0$.

        We will prove $A_t^n \to A_t$ such that $M_t^2 - A_t$ is martingale.

        $\implies A_t = 0$ and $\mathbb{E}M_t^2 = \mathbb{E}M_0^2 = 0.$

        $\implies M_t = 0$ a.s. $\forall~ t \geq 0$.

        By a.s. continuity, $M_t \equiv 0$, a.s.
    }
}

\subsection{It\^o Integral of local martingale}

\thmnn{
    Let $M \in c \mathcal{M}_{loc}$ with $M_0 = 0$. Let $H \in \mathcal{P}$ such that W.P.1
    $$
    \int_0^t H_s ^2 \diff [M]_s < \infty
    $$
    for all $t > 0$. Then $H \cdot M$ is well-defined such that $H\cdot M \in c \mathcal{M}_{0,loc}$.
}
\pf{
    Let $U_n = \inf \{t \colon \abs{ M_t } > n\}$, $V_n = \inf\{t \colon \int_0^t H_s^2 \diff [M]_s > n\}.$

    (Recall $L^2(M) \implies \mathbb{E}\int_0^\infty H_s^2 \diff [M]_s < \infty$).

    $\implies T_n = U_n \wedge V_n$. In this way, $M^{T_n} \in c \mathcal{M}_0^2$ and $H^{T_n} \in L^2(M^{T_n}).$

    So $H^{T_n} \cdot M^{T_n}$ is defined by previous $L^2$ theory.

    $\implies (H\cdot M)^{T_n} = H^{T_n} \cdot M^{T_n}$ ($t < T_n$) is a $L^2$-bounded martingale.

    $\implies H \cdot M \in c \mathcal{M}_{0,loc}$.
}

\defn{semimartingale}{
    We say $X$ is a semimartingale if 
    $$
    X_t = X_0 + M_t + A_t.
    $$
    $M_t$ is local mart, $A_t$ is finite variation.
}

\defnn{
    Let $X$ be a continuous semimartingale. Then $\exists~$unique 
    $$
    \begin{cases}
        M \in c \mathcal{M}_{0,loc}\\
        A \in c F V_0
    \end{cases}
    $$
    such that $X_t = X_0 + M_t + A_t$, $\forall~ t \geq 0$.
}

\pf{
    If $\exists~M'$ and $A'$ such that 
    $$
        X= X_0 + M' + A' = X_0 + M + A.
    $$
    Then $M - M' = A' - A \in c \mathcal{M}_{0,loc} \cap c F V = \{0\}$.
}

We write $M= X^{cm}$ and 
$$
    [X,X] = [M,M] = [X^{cm},X^{cm}].
$$

Similarly, $[X,Y] = [X^{cm},Y^{cm}]$.

For any $H \in lb \mathcal{P}$ define
$$
    H \cdot X =H \cdot M + H \cdot A.
$$
Obviously $H\cdot M \in c \mathcal{M}_{0,loc}$ and $H \cdot A \in FV$.

And 
$$
\begin{aligned}
[H\cdot X, K \cdot Y] &= [H \cdot X^{cm},K\cdot Y^{cm}] \\
&= H \cdot [X^{cm},K\cdot Y^{cm}] \\
&= H\cdot K [X^{cm},Y^{cm}]\\
&= (HK)\cdot [X^{cm},Y^{cm}]=(HK)\cdot [X,Y].
\end{aligned}$$

\subsection{It\^o's Formula} 
\thm{It\^o's Formula}{
    Let $f \colon \mathbb{R}^n \to \mathbb{R}$ be $C^2$. Let 
    $$
    X_t \equiv(X_t^1,\cdots,X_t^n)
    $$
    be a continuous semimartingale.

    Then 
    $$
    f(x_t) - f(x_0) = \int_0^t \nabla f(X_s) \diff X_s + \frac{1}{2}\int_0^t \sum_{i,j} D_{ij}f(X_s)\diff [X^i,X^j]_{s}
    $$
    where $\nabla f(X_s) = (D_1f(X_s),\cdots,D_n f(X_s))$. $D_i f(x) = \frac{\partial}{\partial x_i}f(x)$, $D_{ij} = D_i D_j = D_j D_i.$
}

\thm{Integration by Parts}{
    Let $X,Y$ be continuous semimartingales. Then 
    $$
    X_tY_t - X_0Y_0 = \int_0^t X_s \diff Y_s + \int_0^t Y_s \diff X_s + [X,Y]_t.
    $$
}
\pf{
    Assume IBP holds.
    
    Let $$\mathcal{A} = \left\{f \in C^2 (\mathbb{R}^n) \colon \diff f(X_t) = \sum_i D_i f(X_t) \diff X_t + \frac{1}{2} \sum_{i,j}D_{ij}f(X_t) \diff [X^i,X^j]_t\right\}.$$

    (i). $f,g \in \mathcal{A} \implies f+ g \in \mathcal{A}$.
    (ii). $f,g \in \mathcal{A} \implies fg \in \mathcal{ A}$.

    $$
    \begin{aligned}
    \diff (fg(X_t)) =& \diff (f(X_t)\cdot g(X_t)) \\
    =& f(X_t) \cdot \diff g(X_t) + g(X_t) \diff f(X_t) + \diff [f(X_t),g(X_t)]_t\\
    =& f(X_t)\left( \sum_i D_i g(X_t) \diff X_t + \frac{1}{2} \sum_{i,j} D_{ij}g(X_t) \diff [X^i,X^j]_t \right) \\
    &+ g(X_t)\left( \sum_i D_i f(X_t)\diff X_t + \frac{1}{2}\sum_{i,j}D_{ij}f(X_t) \diff [X^i,X^j]_t \right) \\
    &+ \diff \left[ \int_0^t \sum_{i} D_i f(X_t) \diff X_t, \int_0^t \sum_j D_j g(X_t) \diff X_t^j \right]\\
    =&\diff X_t \left( \sum_i f D_i g(X_t) + \sum_i g D_i f(X_t) \right) \\
    &+ \frac{1}{2} \sum_{i,j} \left( g D_{ij}f (X_t) + f D_{ij}g(X_t) + 2D_i fD_j g(X_t) \right)\diff [X^i,X^j]_t\\
    =& \sum_i D_i(fg)(X_t) \diff X_t + \frac{1}{2}\sum_{i,j}D_{ij}(fg)(X_t)\diff [X^i,X^j]_t\\
    \implies \quad& fg \in \mathcal{A}.
    \end{aligned}
    $$

    (iii). $f(x) = x \in \mathcal{A} \implies \forall~$polynomial $\in \mathcal{A}$.
}

\lemnn{$\forall~ f \in C^2(\mathbb{R}^d)$, $\exists~$polynomials $\{P_n(x)\}_{n \geq 1}$ such that 
    $$
    \begin{cases}
        P_n(x) \to f(x)\\
        D_i P_n(x) \to D_i f(x)\\
        D_{ij}P_n(x) \to D_{ij}f(x)
    \end{cases}
    $$
    uniformly on compacts.
}

Define $U_N = \inf\{t \geq 0\colon \abs{ X_t }+ [M]_{t} > N\}$ for each $N >0$.

Then $U_N \uparrow \infty$. So it suffices to prove It\^o formula holds for $\forall~ t \leq U_N$.

Note that $$P_n(X_t) = P_n(X_0) + \int_{0}^{t}\sum_{i} D_i P_n(X_s) \diff  X_s + \frac{1}{2} \int_{0}^{t} \sum_{i,j} D_{ij}P_n(X_s) \diff [X^i,X^j]_s$$ for $0 \leq t \leq U_N$.

By letting $n \to \infty$, we will prove the limits are 
$$
f(X_t) = f(X_0) + \int_{0}^{t}\sum_{i} D_i f(X_s) \diff X_s + \frac{1}{2} \int_{0}^{t} \sum_{ij}D_{ij}f(X_s) \diff [X^i,X^j]_s.
$$

$$
\begin{aligned}
    I_{1,n}=&\int_{0}^{t} \sum_i D_i P_n(X_s) \diff X_s \\
    =& \int_0^t \sum_i D_i P_n(X_s) \diff M_s + \int_0^t \sum_i D_i P_n(X_s) \diff A_s.
\end{aligned}
$$

Notice that 
$$
    \abs{ \int_0^t \sum_i D_i P_n (X_s) \diff A_s - \int_0^t \sum_i D_i f(X_s) \diff  A_s } \leq d\cdot \sup_{i} \underbrace{\norm{ D_iP_n - D_i f }}_{\text{max on compacts}}\cdot A_t \xlongrightarrow{n \to \infty} 0 \text{ a.s.}
$$

Next,
$$
    \begin{aligned}
    &\mathbb{E} \sup_{s \leq t} \abs{ \int_0^t \sum_i \left( D_i P_n(X_s) - D_i f(X_s) \right) \diff  M_s }^2 \\
    \leq & \mathbb{E}\int_0^t \left( \sum_i D_i P_n(X_s) - D_i f(X_s) \right)^2 \diff [M]_s \\
    \leq & d \cdot \sup_{1 \leq i \leq d}\norm{ D_iP_n - D_i f }^2 \cdot \underbrace{\mathbb{E}[M]_{t \wedge U_N}}_{N} \xlongrightarrow{n \to \infty} 0.
    \end{aligned}
$$

By taking subsequence, $\int_0^t \sum_i D_i P_{n_k}(X_s) \diff  M_s \xlongleftarrow{\text{a.s.}}\int_0^t \sum_i D_i f(X_s) \diff  M_s.$

$$
    \begin{aligned}
        I_{2,n} = \frac{1}{2} \int_0^t \sum_{i,j} D_{ij}P_n(X_s) \diff [X^i,X^j]_s.
    \end{aligned}
$$

Notice that 

$$
    \abs{ I_{2,n} - \frac{1}{2}\int_0^t \sum_{i,j}D_{ij}f(X_s) \diff [X^i,X^j]_s } \leq C \sum_{i,j} \norm{ D_{ij}P_n - D_{ij}f }\cdot \underbrace{\sum_{i,j}[X^i,X^j]_t}_{\leq C_N} \xlongrightarrow{n \to \infty} \to 0.
$$

\lemnn{
    Let $M$ be a bounded continuous martingale with $M_0 = 0$. Let $V$ be continuous $FV_0$
    $$
    \begin{aligned}
        &M_t^2 = \int_0^t 2M_s \diff M_s + [M]_t ,\\
        &M_t V_t = \int_0^t M_s \diff V_s + \int_0^t V_s \diff  M_s.
    \end{aligned}
    $$
}
\pf{
    $$
    \begin{aligned}
        M_t^2 &= \left[ \sum_{k=1}^{2^n}M(t_{k}^{n}) - M(t_{k-1}^{n}) \right]^2\\
        &= \sum_{k=1}^{2^n}\left( M(t_{k}^{n}) - M(t_{k-1}^{n}) \right)^2 + 2\sum_{k=1}^{2^n}\left[ M(t_{k}^{n})M(t_{k-1}^{n}) - M(t_{k-1}^{n})^2 \right]\\
        &=M(t_{2^n}^{n})^2 - M(t_{0}^{n})^2 = M(t)^2
    \end{aligned}
    $$
    where $t_{k}^{n} = \frac{k}{2^n} t$.

    Note that $\sum_{k=1}^{2^n}\left( M(t_{k}^{n}) - M(t_{k-1}^{n}) \right)^2 \to [M]_t$ and 
    $$
    2\sum_{k=1}^{2^n}\left[ M(t_{k}^{n})M(t_{k-1}^{n}) - M(t_{k-1}^{n})^2 \right] = 2 \sum_{k=1}^{2^n}M(t_{k-1}^{n})(M(t_{k}^{n}) - M(t_{k-1}^{n}))\to 2\int_0^t M_s \diff M_s.
    $$

    $$
    \begin{aligned}
        M_t V_t &= M_{t_{2^n}^n}V_{t_{2^n}^n} - M_{t_{0}^n}V_{t_0^n} \\
        &= \sum_{k=1}^{2^n}\left[ M(t_{k}^{n})V(t_{k}^{n}) - M(t_{k-1}^{n})V(t_{k-1}^{n}) \right] \\
        &= \sum_{k=1}^{2^n} M(t_{k}^{n})(V(t_{k}^{n}) - V(t_{k-1}^{n})) + \sum_{k=1}^n V(t_{k-1}^{n})\left( M(t_{k}^{n}) - M(t_{k-1}^{n}) \right) \\
        &\to \int_{0}^{t}M_s \diff V_s + \int_0^t V_s \diff M_s.
    \end{aligned}
    $$
}

\thm{L\'evy Thm}{
    Let $\{B_t^i\}_{t \geq 0} \in c \mathcal{M}_{0,loc}$ for $i=1,2,\cdots,n$. If $\forall~i,j$, we have 
    $$ 
    [B^i,B^j]_t = \begin{cases}
        t, & i = j,\\
        0, & i \neq j.
    \end{cases}
    $$
    
    Then $B_t = (B_t^1,\cdots,B_t^n)$ is a Brownian motion on $\mathbb{R}^n$ with $\mathcal{N}(0,I_n).$
}

\pf{
    $\forall~ s < t,$ we will prove 
    $$
    \mathbb{E}\left[ e^{i \theta (B_t - B_s)} \middle| \mathcal{F}_s \right] = e^{-\frac{1}{2}(t-s) \abs{ \theta }^2}
    $$
    for all $\theta \in \mathbb{R}^n$.

    Define $f(x,t) = e^{i\theta x + \frac{1}{2}\abs{ \theta }^2 t }$, $\mathbb{R}^{n+1} \to \mathbb{C}$.

    Apply It\^o's formula to $f(B_t,t) \equiv f(X_t).$

    $X_t = (B_t^1,B_t^2,\cdots,B_t^n,t).$

    $$
    \begin{aligned}
    f(X_t) =& f(x_0) + \int_0^t \sum_{i=1}^n (D_i f(X_s \diff B_s^i)) + \int_0^1 D_{n+1}f(X_s)ds \\
    &+ \frac{1}{2}\int_0^t \sum_{1 \leq i,j \leq n}D_{ij}f(X_s) \diff [B^i,B^j]_s \\
    &+ \frac{1}{2}\int_0^t \sum_{1 \leq k \leq n}D_{k,n+1}f(X_s) \diff [B^k,s]_s
    \end{aligned}
    $$
    $$
    \begin{aligned}
        \implies f(B_t,t)= &f(B_0,0) + \int_0^t \sum_i D_i f(B_s,s)\diff B_s^i \\
        &+ \int_0^t \left( e^{i \theta B_s + \frac{1}{2}\abs{ \theta }^2 s} \cdot \frac{1}{2} \abs{ \theta }^2 \right) \diff s \\
        &+ \int_0^t \frac{1}{2}\left( \sum_{j=1}^n e^{i \theta B_s + \frac{1}{2}\abs{ \theta }^2 s }(i \theta_j)^2 \right) \diff s.
    \end{aligned}
    $$

    We conclude that $f(B_t,t) = f(B_0,0) + c \mathcal{M}_{0,loc}$
    $$
    \implies e\underbrace{^{i\theta B_t + \frac{1}{2} \abs{ \theta }^2 t}}_{M_t^\theta \in c \mathcal{M}_{0,loc}} = 1 + c \mathcal{M}_{0,loc}
    $$
    Since $\abs{ M_t^\theta } = e^{\frac{1}{2} \abs{ \theta }^2 t} < \infty$ and 
    $$
    \mathbb{E}(M_t^{T_n} \mid \mathcal{F}_s) = M_s^{T_n}
    $$
    where $T_n$ reduces $M$.

    We may take $n \to \infty$ to get $\mathbb{E}(M_t \mid \mathcal{F}_s) = M_s$.

    $$
    \implies \mathbb{E}(e^{i \theta B_t + \frac{1}{2}\abs{ \theta }^2 t} \mid \mathcal{F}_s) = e^{i\theta B_s + \frac{1}{2} \abs{ \theta }^2 s}.
    $$

    $$
    \implies \mathbb{E}\left( e^{i \theta(B_t - B_s)} \mid \mathcal{F}_s \right) = e^{\frac{1}{2}\abs{ \theta }^2 (t-s)}.
    $$

    $$
    \implies B_t - B_s \sim \mathcal{N}(0,t-s)
    $$

    Q independent of $\mathcal{F}_s$.
}

\thm{Dubins and Schwarz}{
    let $M \in c \mathcal{M}_{0,loc}$ such that $[M]_t \uparrow \infty$ as $t \to \infty$. Then $M_t = B([M]_t)$ where $B$ is some Brownian motion.

    Let $\tau_t = \inf\{u \geq 0 \colon [M]_u > t\}$. Set $\mathcal{G}_t \equiv \mathcal{F}(\tau_t)$. Then $B_t \equiv M(\tau_t)$ is a B.M. with respect to $\mathcal{G}_t.$
}

\pf{
    (i) $M(\tau_t) \in \mathcal{F}_{\tau_t} \implies B_t \in \mathcal{G}_t$.
    
    (ii) $t\mapsto B_t$ is continuous.

    $$
    \begin{cases}
        \tau_{a-} = b, \\
        \tau_{a+} = c.
    \end{cases}
    $$

    We only need to prove $M(b) = M(c)$.

    Assume $M$ is bounded $\forall~ q \in \mathbb{Q}_+$, define 
    $$
    S_q \equiv \inf\{t > q \colon [M]_t > [M]_q\}.
    $$

    We will show $M_t \equiv M_q$ for all $t \in [q,S_q)$
    
    Note $S_q \geq q$, by OST,
    $$\mathbb{E}\left( M_{S_q}^2 - [M]_{S_q} \mid \mathcal{F}_q \right) = M_q^2 - [M]_q.
    $$

    Because $[M]_{S_q} = [M]_q$,
    $$
    \implies \mathbb{E}[M_{S_q}^2 - M_q^2 \mid \mathcal{F}_q] = 0.
    $$
    $$
    \implies \mathbb{E}((M_{S_q} - M_q)^2 \mid \mathcal{F}_q) = 0.
    $$
    $$
    \implies \mathbb{E}(M_{S_q} - M_q)^2 = 0.
    $$
    $$
    \implies M_{S_q} = M_q.
    $$

    (iii) $B_t$ and $B_t^2 - t \in c \mathcal{M}_{0,loc} \implies [B]_t = t \implies$ B.M.

    Recall $B_t = M(\tau_t)$. Define 
    $$
    \begin{aligned}
    &T_n = \inf \{ t \geq 0 \colon \abs{ M_t } > n\},\\
    &U_n = [M]_{T_n}.    
    \end{aligned}
    $$

    We will check $\tau_{t \wedge U_n} = T_n \wedge \tau_t.$
    $$
    t \wedge U_n \leq t \implies \tau_{t \wedge U_n } \leq \tau_t.
    $$
    $$
    \tau_{t\wedge U_n} =\inf \{
        u\colon [M]_u > t \wedge U_n = t \wedge[M]_{T_n}
    \}.
    $$

    \textbf{Claim:} $\tau_t \wedge T_n = \tau_{t \wedge U_n}$
    
    \textbf{case 1.} $t \leq U_{n}$ we have
    $$
        \tau_{t \wedge U_{n}} = \inf \{u\colon [M]_{u}> t\} = \tau_{t},
    $$
    and
    $$
    t \leq U_{n} = [M]_{T_{n}} \implies T_{n} \geq \tau_{t} \implies \tau_{t} \wedge T_{n} = \tau_{t} = \tau_{t \wedge U_{n}}.
    $$

    \textbf{case 2.} $t > U_{n}$ we have
    $$
    \tau_{t \wedge U_{n}} = \tau_{U_{n}} = \inf \{u\colon [M]_{u} > U_{n}\}
    $$
    and
    $$
    t > U_{n} > [M]_{T_{n}} \implies T_{n} < \tau_{t} \implies \tau_{t} \wedge T_{n} = T_{n}.
    $$
    It suffices to prove $T_{n} = \tau_{[M]_{T_{n}}}= \tau_{U_{n}}$.

    One hand, $\tau_{U_{n}} \geq T_{n}$ can be proved by def of $\tau_{U_{n}}$.

    On the other hand, $T_{n} = \inf \{t \colon \abs{ M_{t} } > n\}$. Take $\forall~\varepsilon>0$, $u_{\varepsilon} \in (T_{n} , T_{n} + \varepsilon)$ s.t. $\abs{ M_{u_{\varepsilon}} }>n .$

    Hence $M_{u_{\varepsilon}} \neq M_{T_{n}}$. Suppose $[M]_{u_{\varepsilon}} = [M]_{T_{n}}$, then $M$ is constant by the property of continuous local martingale. Contradiction. So $[M]_{u_{\varepsilon}} > [M]_{T_{n}} = U_{n} \implies \tau_{U_{n}} \leq u_{\varepsilon}$. Let $\varepsilon \downarrow 0$ to get $\tau_{U_{n}} \leq T_{n}$.

    Then $B_{t \wedge U_n} = M(\tau_{t \wedge U_n}) = M(T_n \wedge \tau _t)$, $B_t^{U_n} = M_{\tau_t}^{T_n}.$

    By OST 
    $$
    \mathbb{E}\left( B_t^{U_n} \middle| \mathcal{F}_{\tau_s} \right) = \mathbb{E}\left( M_{\tau_t}^{T_n} \middle| \mathcal{F}_{\tau_s} \right) = M_{\tau_s}^{T_n} = B_s^{U_n}.
    $$

    Then $B^{U_n}$ is a martingale w.r.t. $\mathcal{G}_t = \mathcal{F}_{\tau_t}$ + (2) $U_n$ is $\mathcal{G}_t$-stopping time $\implies B \in c \mathcal{M}_{0,loc}$

    \textbf{pf of (2)}

    $$
    \{[M]_{T_n} \leq t\}=\{U_n \leq t\} \in \mathcal{F}_{\tau_t}.
    $$

    $$
    \{T_n \leq \tau_t\} \in \mathcal{F}_{\tau_t}.
    $$

    $\left( M^2 - [M] \right)^{T_n}$ is martingale. Then 
    $$
    \mathbb{E}\left[ (M_{\tau_t}^{T_n})^2 - [M]_{\tau_t}^{T_n} \middle|  \mathcal{F}_{\tau_s}\right]= \left( M_{\tau_s}^{T_n} \right)^2 - [M]_{\tau_s}^{T_n}
    $$  
    $$
    \implies \mathbb{E}\left[ 
        (B_t^{U_n})^2 - t \wedge U_n \middle| \mathcal{F}_s 
     \right] = (B_s^{U_n})^2 - s \wedge U_n.
    $$  
    $
    \implies (B_t^2 - t)^{U_n}
    $ is a martingale $\implies B_t^2 - t \in \mathcal{M}_{loc} \implies[B]_t = t.$
}   

If $[M]_{\infty} < \infty$, we get $\tau_t = \infty$, $\forall~ t \geq [M]_{\infty}$.
$$
B_t = M(\infty), \quad \forall~ t \geq [M]_{\infty}.
$$

So this def does not give a B.M.

Define 
$$B_t = 
\begin{cases}
    M(\tau_t), & t < [M]_{\infty},\\
    M(\infty)+W(t) - W([M]_{\infty}), & t \geq [M]_{\infty},
\end{cases}
$$
where $W$ is a standard B.M. independent of $M$.

Or 
$$B_t
\begin{cases}
    M(\tau_t),& t <[M]_{\infty}\\
    M(\infty)+W(t-[M]_{\infty}), & t \geq[M]_{\infty}.
\end{cases}
$$
Then $B([M]_{t})= M_t.$

\lemnn{
    On $\{[M]_{\infty
} < \infty\}$, the limit $M_\infty \equiv \lim_{t \uparrow \infty} M_t$ exists a.s.
}
\pf{
    Define $S_K = \inf \{t \geq 0 \colon [M]_t > K\}$. Then on $\{[M]_{\infty} < \infty\}$, $S_K = \infty$, $\forall~ K >[M]_\infty$

    $M^{S_K}$ is a $L^2$-bounded martingale.

    $\implies \lim_{t \uparrow \infty} M_t^{S_K}$ exists a.s.

    Let $K \uparrow \infty$ to get $\lim_{t\uparrow \infty}M_t$ exists a.s.

    By an enlargement of the prob. space $(\Omega. \mathcal{F},\mathsf{P})$
    $$
    \begin{aligned}
    \implies &\tilde{\Omega} = \Omega \times \Omega^*\\
        & \tilde{\mathcal{F}}_t = \mathcal{F}_t \times \mathcal{F}_{t}^* \\
        & \tilde{\mathsf{P}} = \mathsf{P} \times \mathsf{P}^*
    \end{aligned}
    $$
    where on $(\Omega^*, \mathcal{F}^*,\mathcal{F}_{t}^*, \mathsf{P}^*)$ we have $W^*$ is B.M. independent of $M$.
}

\thm{Dubins-Schwarz (general version)}{
    Let $M \in c \mathcal{M}_{0,loc}$. Then $\exists~$B.M. B s.t. $M_t = B([M]_t)$ for all $t \geq 0$.
}

\thm{Dol\'eans' Thm}{
    Let $X$ be a semimartingale with $X_0 = 0$. Suppose $Z_0 \in \mathcal{F}_0$. Then $\exists~$a unique continuous semimartingale $Z$ s.t. 
    $$
    Z_t = Z_0 + \int_{0}^{t}Z_s \diff X_s.
    $$
    The solution $Z_t = Z_0 e^{X_t - \frac{1}{2}[X]_t}$. 

    \textbf{Notation:} $\mathcal{E}(X_t) = e^{X_t - \frac{1}{2}[X]_t}$ is called the exponential semimartingale of $X$.
}

\pf{
    $\mathcal{E}(X_t) = e^{X_t - \frac{1}{2}[X]_t} = f(X_t, [X]_t)$ with $f(x,y) = e^{x - \frac{1}{2}y}.$

    By apply It\^o's formula, 
    $$
    \diff f(X_t,[X]_t) = \partial_x f(X_t,[X]_t) \diff X_t + \partial_y f(X_t, [X]_t) \diff [X]_t + \frac{1}{2}\left( \partial_{xx}f(X_t,[X]_t) \diff [X]_t \right).
    $$
    Because 
    $$
    \begin{cases}
        \diff X_t \diff X_t = \diff [X]_t\\
        \diff X_t \diff [X]_t = 0\\
        \diff [X]_t \diff [X]_t = 0.
    \end{cases}
    $$
    Hence 
    $$
    \diff \mathcal{E}(X_t) = \mathcal{E}(X_t) \diff X_t.
    $$
    So $\mathcal{E}(X)$ is a solution to 
    $$
    \diff Z_t = Z_t \diff X_t.
    $$

    \textbf{(Uniqueness)} Assume $Z$ is another solution. We will prove 
    $$
        Z_t = Z_0 \mathcal{E}(X_t).
    $$

    We now have $\diff Z_t = Z_t \diff X_t$. 
    
    Define 
    $$
    Y_t = e^{- X_t + \frac{1}{2}[X]_t} = (\mathcal{E}(X)_t)^{-1}.
    $$

    It suffices to prove $\diff (Z_t Y_t) = 0$ which implies 
    $$
    Z_t Y_t = Z_0 Y_0 = Z_0, \quad \forall~t.
    $$
    
    We have
    $$
    \diff (Z_t Y_t) = Z_t \diff Y_t + Y_t \diff Z_t + \diff [Y,Z]_t.
    $$

    Notice that by It\^o's formula,
    $$
    \diff Y_t = (-1)Y_t \diff X_t + \frac{1}{2} Y_t \diff [X]_t + \frac{1}{2}(-1)^2 Y_t \diff [X]_t = -Y_t \diff X_t + Y_t \diff [X]_t.
    $$

    $$
    \implies \diff (Z_t Y_t) = Z_t(-Y_t \diff X_t + Y_t \diff [X]_t) + Y_t(Z_t \diff X_t) + Y_t(Z_t \diff X_t) + (-Y_t Z_t) \diff [X]_t = 0.
    $$
}

Let $M \in c \mathcal{M}_{0,loc}$. Then $\mathcal{E}(\theta M)_t = e^{\theta M_t - \frac{1}{2}\theta^2[M]_t}$ is in $c \mathcal{M}_{loc}$.

By Fatou's lemma, we get 
$$
\begin{aligned}
    \mathbb{E}(\mathcal{E}(\theta M)_t \mid \mathcal{F}_s) &= \mathbb{E}(\lim_{n\to \infty} \mathcal{E}(\theta M)_{t \wedge T_n} \mid \mathcal{F}_s) \\
    &\leq \liminf_{n \to \infty} \mathbb{E}(\mathcal{E}(\theta M)_{t}^{T_n} \mid \mathcal{F}_s)\\
    & = \liminf_{n \to \infty} \mathcal{E}(\theta M)_{s}^{T_n} \\
    &= \mathcal{E}(\theta M)_s.
\end{aligned}
$$

Hence $\mathcal{E}(X_t)$ is a supermartingale and $\mathbb{E}(\mathcal{E}(M)_t) \leq 1$.

\thmnn{
    Suppose $\forall~ t > 0$, $\exists~ K_t > 0$ s.t.
    $$
    [M]_t < K_t \quad \text{a.s.}
    $$

    If for some $t > 0$, $\exists~K_t >0$ s.t.
    $$
    [M]_t < K_t \quad \text{a.s.}
    $$
    
    Then for this $t > 0$,
    $$
    \mathsf{P}\left( \max_{s \leq t}M_s > y \right) \leq e^{-\frac{y^2}{2K_t}}, \quad \forall~ y > 0.
    $$
    and $\mathcal{E}(\theta M)$ is a true martingale.
}

\pf{
    By using $\mathcal{E}(\theta M)$ is a supermartingale, by Doob's inequality,
    $$
    \begin{aligned}
    \mathsf{P}\left( \max_{s \leq t} M_s > y \right) &= \mathsf{P}\left( \max_{s \leq t} \mathcal{E}(\theta M)_s > e^{\theta y - \frac{1}{2} \theta^2 K_t} \right)\\
    & \leq \frac{\mathbb{E}(\mathcal{E}(\theta M)_t)}{e^{\theta y -\frac{1}{2} \theta^2 K_t}}\\
    & \leq e^{-\theta y + \frac{1}{2} \theta^2 K_t}.
    \end{aligned}
    $$
    Pick $\theta = \frac{y}{K_t}$ to get $e^{-\frac{y^2}{2K_t}}$.

    Next, define $Z_t = \max_{s \leq t} M_s$. Then $\forall~ \theta > 0$, 
    $$
        \mathbb{E}e^{\theta Z_t} = \int_{\mathbb{R}}e^{\theta y} \mathsf{P}(Z_t \in \diff y).
    $$

    Notice that $\forall~ z \in \mathbb{R}$, 
    $$
        e^{\theta Z} = 1 + \int_{0}^{Z} \theta e^{\theta y} \diff y.
    $$
    $$
    \begin{aligned}
    \implies 
        \mathbb{E}(e^{\theta Z_t}) &= 1 + \mathbb{E}\left( \int_0^{Z_t} \theta e^{\theta y} \diff y \right) \\
        &= 1 + \mathbb{E}\left( \int_0^\infty \mathsf{1}_{y < Z_t} \theta e^{\theta y} \diff  y \right)\\
        & = 1+ \int_0^\infty \theta e^{\theta y} \diff  y  \mathsf{P}(Z_t > y)\\
        & \leq 1 + \int_0^\infty \theta e^{\theta y} e^{-\frac{y^2}{2K_t}} \diff y \\
        &< \infty.
    \end{aligned}
    $$

    By $\mathcal{E}(\theta M)_t$ is a $\mathcal{M}_{loc}$, we get $\exists~T_n \uparrow \infty$ s.t. $\mathcal{E}(\theta M)_{t}^{T_n}$ is a martingale.

    $$
        \implies \forall~ s < t,\quad \mathbb{E}[\mathcal{E}(\theta M)_{t}^{T_n} \mid \mathcal{F}_s] = \mathcal{E}(\theta M)_{s}^{T_n}.
    $$

    By DCT for $\mathcal{E}(\theta M)_{t}^{T_n} \leq Z_t$ and $\mathbb{E} Z_t < \infty$, we get 
    $$
    \lim_{n \to \infty} \mathbb{E}(\mathcal{E}(\theta M)_{t}^{T_n} \mid \mathcal{F}_s) = \mathbb{E}(\mathcal{E}(\theta M)_t \mid \mathcal{F}_s).
    $$
}

\cor{
    Let $B$ be a B.M. in $\mathbb{R}^n$ and let $H$ be a bounded previsible process. Then 
    $$
    e^{\int H \diff B - \frac{1}{2}\int \abs{ H_s }^2 \diff s}
    $$
    is a true martingale.
}

\rmkb{$[M]_t = \int_0^t \abs{ H_s }^2 \diff  s < K_t$.}

\cor{
    Let $M \in c \mathcal{M}_{0,loc}$. Then $\forall~ t > 0$, $y> 0$, $K >0$,
    $$
    \mathsf{P}\left( \max_{s \leq t} M_s \geq y, [M]_{t} \leq K \right) \leq e^{-\frac{y^2}{2K}}.
    $$
    $$
    \implies \mathsf{P}\left( \max_{s \leq t}\abs{ M_s } \geq y, [M]_t \leq K \right) \leq 2 e^{-\frac{y^2}{2K}}.
    $$
}

\pf{
    Let $T =\inf \{u \geq 0 \colon [M]_{u} \geq K\}$. Then 
    $$
    \mathsf{P}\left( \max_{s \leq t}M_s \geq y, [M]_{t} \leq K \right) \leq \mathsf{P}\left( \max_{s \leq t} M_{s \wedge T} \geq y \right) \leq e^{- \frac{y^2}{2K}}
    $$  
    since $[M^T]_t \leq K$.
}

\thm{Novikov's Thm}{
    If $\forall~t > 0$, $\mathbb{E}(e^{\frac{1}{2} [M]_t} ) < \infty$ then $\mathcal{E}(X)_t$ is a true martingale.
}

\pf{
    See P198 of [KS].
}

\subsection{Burkholder-Davis-Gundy inequalities}

\defnn{
    We say a function $F: \mathbb{R}^+ \to \mathbb{R}^+$ is \underline{moderate} if continuous and increasing, $F(x) = 0 \Leftrightarrow x=0$, and $\forall~\alpha > 1$,
    $$
        \sup_{x > 0} \frac{F(\alpha x)}{F(x)} < \infty.
    $$
}
\ex{}{
    $F(x) = x^{p}$, $\forall~ p > 0$.
    $F(x) = x^p \left( \log x \vee 1 \right)$, $\forall~ p \geq 0$.
}

\thm{BDG inequality}{
    Let $F$ be moderate. Then $\exists~$universal constants $c_F \leq C_F < \infty$ s.t. $\forall~ M \in c \mathcal{M}_{0,loc}$,
    $$
    c_{F} \mathbb{E}\left( F([M]_{\infty}^{\frac{1}{2}}) \right) \leq \mathbb{E} F(M_\infty^*) \leq C_F \mathbb{E}\left( F([M]_{\infty}^{\frac{1}{2}}) \right).
    $$

    $M_t^* = \sup_{s \leq t}\abs{ M_s }.$ Fix $t > 0$. Let $\tilde{M}_{s}^{t} = M_{s \wedge t}$. Then 
    $$
    c_{F} \mathbb{E}\left( F([M]_{t}^{\frac{1}{2}}) \right) \leq \mathbb{E}\left( F(M_t^*) \right) \leq C_F \mathbb{E} F([M]_{t}^{\frac{1}{2}}).
    $$
}

\ex{}{
    If $F(x) = x^p$, then 
    $$
    c_p \mathbb{E}[M]_{t}^{\frac{p}{2}} \leq \mathbb{E} \abs{ M_t^* }^p \leq C_p \mathbb{E}[M]_t^{\frac{p}{2}}.
    $$
}

\lemp{``good'' $\lambda$ inequality}{
    Let $X$ and $Y$ be nonegative R.V.s. Suppose $\exists~\beta > 1$ s.t. $\forall~ \lambda > 0$, $\delta >0$,
    $$
    \mathsf{P}(X > \beta \lambda, Y \leq \delta \lambda) \leq \psi(\delta) \mathsf{P}(X > \lambda)
    $$
    where $\psi(\delta) \downarrow 0$ as $\delta \downarrow 0$.

    Then for each moderate $F$, $\exists~C=C(\beta,\psi, F) > 0$ s.t. 
    $$
    \mathbb{E}[F(X)] \leq c \mathbb{E}[F(Y)].
    $$
}{
    Assume good $\lambda$ inequality. Then it holds by replacing $X$ by $X \wedge a$ for any $a>0$.

    If $a > \beta \lambda > \lambda$, then
    $$
    \begin{aligned}
    &X \wedge a > \beta \lambda \Leftrightarrow X > \beta \lambda\\
    & X \wedge a > \lambda \Leftrightarrow X > \lambda.
    \end{aligned}
    $$

    Hence we may assume $\mathbb{E}F(X) < \infty$.

    [ $\mathbb{E}F(X\wedge a) \leq \mathbb{E} F(Y)$ for all $a >0$. Then let $a \uparrow \infty$.]

    Notice that $\sup_{x>0} \frac{F(x)}{F(\frac{x}{\beta})}$.

    Let $\gamma > 0$ s.t. $F(x) \leq \frac{1}{\gamma} F(\frac{x}{\beta})$ for all $x > 0$.

    Now 
    $$
    \int F(\diff \lambda) \psi(\delta) \mathsf{P}(X > \lambda) \geq \int F(\diff \lambda) \mathsf{P}(X > \beta \lambda, Y \leq \delta \lambda).
    $$
    $$
    LHS = \psi(\delta) \mathbb{E}\left( \int_0^X F(\diff \lambda) \right) = \psi(\delta)\mathbb{E}(F(X)).
    $$
    $$
    \begin{aligned}
    RHS &= \int F(\diff \lambda) \mathbb{E}(\mathsf{1}_{\frac{Y}{\delta} \leq \lambda < \frac{X}{\beta}}) \\
    &= \mathbb{E}\left( \int_{\frac{Y}{\delta}}^{\frac{X}{\beta}} F(\diff \lambda) \right) \\
    &= \mathbb{E}\left( F\left(\frac{X}{\beta}\right) - F\left(\frac{Y}{\delta}\right) \right)\\
    & \geq \mathbb{E}\left( \gamma F(X) \right) - \mathbb{E}\left( F\left( \frac{Y}{\delta} \right) \right).
    \end{aligned}
    $$

    Let $\delta > 0$ be small sucht that 
    $$
    \psi(\delta) < \frac{\gamma}{2} \implies \gamma -\psi(\delta) \geq \frac{\gamma}{2}.
    $$

    Hence 
    $$
    \mathbb{E}\left( F(X) \right)\left( \psi(\delta) - \gamma \right) \geq - \mathbb{E}\left( F\left( \frac{Y}{\delta} \right) \right).
    $$

    $$ \implies
    \mathbb{E}\left( F\left( \frac{Y}{\delta} \right) \right) \geq \left( \gamma - \psi(\delta) \right) \mathbb{E}F(X) \geq \frac{\lambda}{2} \mathbb{E}F(X).
    $$

    Using $\displaystyle\sup_{y > 0}\frac{F(y/\delta)}{F(y)} < \infty$, $\implies \exists~\mu > 0$ s.t. $F(y/\delta) \leq \mu F(y)$ for all $y > 0$ 
    $$\implies \mathbb{E}(F(Y/\delta)) \leq \mu \mathbb{E} F(Y)$$
    $$
    \implies \frac{\gamma}{2}\mathbb{E}F(X) \leq \mu \mathbb{E}F(Y).
    $$
}

\pf{
    of BDF inequality.

    $$\mathsf{P}\left(M_\infty^* > \beta \lambda , [M]_{\infty}^{\frac{1}{2}} \leq \delta \lambda\right) \leq e^{-\frac{(\beta \lambda)^2}{2 (\delta \lambda)^2}} = e^{-\frac{\beta^2}{2 \delta^2}}.$$

    Define $\tau = \inf \{u \geq 0 \colon \abs{ M_u }> \lambda\}$. 

    Then 
    $$\begin{aligned}
        \mathsf{P}\left( M_\infty^* > \beta \lambda, [M]_{\infty}^{\frac{1}{2}} \leq \delta \lambda \right) &= \mathsf{P}\left( M_\infty^* > \beta \lambda, [M]_{\infty}^{\frac{1}{2}} \leq \delta \lambda , \tau < \infty\right)\\
        &= \mathbb{E}\left[ \mathsf{1}_{\tau < \infty} \mathsf{P}\left( M_\infty^* > \beta \lambda, [M]_{\infty}^{\frac{1}{2}} \leq \delta \lambda \,\middle|\, \mathcal{F}_{\tau} \right) \right]
    \end{aligned}
    $$

    Define $N_t \equiv (M_{t + \tau} - M_{\tau})^2 - ([M]_{t + \tau} - [M]_{\tau})$.

    Then $N_t \in \mathcal{F}_{t + \tau}$ is a continuous local martingale
    $$
    \mathbb{E}\left[ M_{t+\tau}^2 - [M]_{t + \tau} \mid \mathcal{F}_{s + \tau} \right] = M_{s+\tau}^2 - [M]_{s+\tau}.
    $$

    $ \implies N_0 = 0$ and on $\{M_{\infty}^* > \beta \lambda, [M]_{\infty}^{\frac{1}{2}} \leq \delta \lambda\}$ we have $\exists~\tilde{t} > 0$ s.t.
    $$
    M_{\tilde{t} + \tau} \geq \beta \lambda \implies N_{\tilde{t}} = \left( M_{\tilde{t}+\tau} \pm \lambda \right)^2 - \left( [M]_{\tilde{t} + \tau} - [M]_{\tau} \right) \geq \left[ (\beta - 1) \lambda \right]^2 - (\delta \lambda)^2.
    $$

    Moreover $N_t \geq - \delta^2 \lambda^2$, $\forall~ t \geq 0$. 

    Hence $N_t$ reaches $[(\beta-1)\lambda]^2 - (\delta \lambda)^2$ before $(-\delta^2 \lambda^2)$.

    By (Gambler's Ruin $\mathsf{P}_{k}(\tau_0 < \tau_N) = \frac{N-k}{N}$)
    $$ 
    \mathsf{P}(\tau_b < \tau_a) \leq \frac{-a}{b-a}
    $$
    $$
    \implies \mathsf{P}\left( M_\infty^* > \beta \lambda, [M]_{\infty}^{\frac{1}{2}} \leq \delta \lambda \mid \mathcal{F}_\tau \right) \leq \mathsf{P}\left( \tau_{[(\beta - 1)\lambda]^2 - (\delta \lambda)^2} < \tau_{-(\delta \lambda)^2} \right) \leq \frac{(\delta \lambda)^2}{[(\beta -1)\lambda]^2} = \frac{\delta^2}{(\beta -1)^2}.
    $$

    Here we take $0 < \delta < \beta -1$. Therefore
    $$
    \mathsf{P}\left( M_\infty^* > \beta \lambda, [M]_{\infty}^{\frac{1}{2}} \leq \delta \lambda  \right) \leq \delta^2 \frac{1}{(\beta - 1)^2} \mathsf{P}(\tau < \infty) = \delta^2 \frac{1}{(\beta-1)^2} \mathsf{P}\left( M_\infty^* > \lambda \right).
    $$

    Similarly, 
    $$
    \mathsf{P}\left( [M]_{\infty}^{\frac{1}{2}} > \beta \lambda, M_{\infty}^*  \leq \delta \lambda\right).
    $$

    $$
    \tilde{N}_t = [M]_{t+\tilde{\tau}} - [M]_{\tilde{\tau}} - (M_{\tilde{\tau} + t}-M_{\tilde{\tau}})^2
    $$
    with $\tilde{\tau} = \inf\{u \geq 0 \colon [M]_{u}^{\frac{1}{2}} > \lambda\}$.

    $$
    \implies \mathsf{P}\left( [M]_{\infty}^{\frac{1}{2}} > \beta \lambda, M_{\infty}^*  \leq \delta \lambda, \tilde{\tau} < \infty \right) = \mathbb{E}\left[ \mathsf{P}(\dots \mid \mathcal{F}_{\tilde{\tau}})\mathsf{1}_{\tilde{\tau} < \infty} \right].
    $$

    Then $\tilde{N}_t$ reaches $(\beta \lambda)^2 - \lambda^2 - (\delta \lambda)^2$ before $-(\delta \lambda)^2$.

}

\section{Semimartingale local time}
For a Brownian motion $(B_t)_{t \geq 0}$. We consider
$$
l_t^0 = \int_0^t \delta_0 (B_s) \diff s,
$$
where $\delta_0$ is the Dirac delta function such that $\forall~$function $f$,
$$
\int f(x) \delta_0(x) \diff x = f(0).
$$

$$
\mu([a,b]) = l_b^0 - l_a^0,
$$
$\mathrm{supp}~(\mu)=\{t \geq 0 \colon B_t = 0\}$, level set $=\{t \geq 0 \colon B_t = x\}$, using local time 
$$
l_t^x = \int_0^t \delta_x (B_s) \diff s.
$$

\thm{Tanaka's formula}{
    Let $X$ be a continuous semimartingale. Then $\exists~$continuous increasing adapted process $\{l_t, t \geq0\}$
    $$
    \abs{ X_t } - \abs{ X_0 } = \int_0^t \mathrm{sgn}(X_s) \diff X_s + l_t,
    $$
    where 
    $$
    \mathrm{sgn}(x) = \begin{cases}
        1, & x >0\\
        0, & x = 0\\
        -1,&x< 0.
    \end{cases}
    $$
    Then $\{l_t , t \geq 0\}$ is called the local time of $X$ at 0, and 
    $$\int_0^t \mathsf{1}_{\{X_s \neq 0\}} \diff l_s = 0.$$
}

\rmkb{
    \begin{itemize}
        \item[(1)] $f(x) = \abs{ x } \implies f'(x) = \mathrm{sgn}(x)$, $\forall~ x \neq 0 \implies f''(x) = 2 \delta(x).$ (distributional derivative)
        \item[(2)] $\abs{ X_t - x} - \abs{ X_0 - x } = \int_0^t \mathrm{sgn}(X_0 - x) \diff X_s + l_t^x.$
        \item[(3)] Note that $\abs{ X_t } = X_t^+ + X_t^-$, $X_t = X_t^+ - X_t^- $
        $$\implies 2X_t^+ = \abs{ X_t } + X_t = \int_0^t \mathsf{1}_{\{X_s > 0\}} \diff X_s - \int_0^t \mathsf{1}_{\{X_s \leq 0\}} \diff X_s + l_t + \int_0^t 1 \diff X_s + X_0 + \abs{ X_0 }$$
        $$
        \implies 2X_t^+ = 2 \int_0^t \mathsf{1}_{\{X_s > 0\}} \diff X_s + l_t + 2X_0^+.
        $$
        $$
        \implies X_t^+ - X_0^+ = \int_0^t \mathsf{1}_{\{X_s > 0\}} \diff X_s + \frac{1}{2} l_t.
        $$

        Similarly, $2X_t^- = \abs{ X_t } - X_t$
        $$
        \implies X_t^- - X_0^- = -\int_0^t \mathsf{1}_{\{X_s \leq 0\}} \diff X_s + \frac{1}{2} l_t.
        $$
    \end{itemize}
}

\pf{
    Apply It\^o's formula to the function $f(x) = \abs{ x }.$
    $$
    \begin{aligned}
    \implies \abs{ X_t }&= \abs{ X_0 } + \int_0^t f'(X_s) \diff X_s + \frac{1}{2}\int_0^t f''(X_s) \diff [X]_s\\
    &= \abs{ X_0 } + \int_0^t \mathrm{sgn}(X_s) \diff X_s + \frac{1}{2} \underbrace{\int_0^t 2 \delta(X_s) \diff [X]_s}_{=l_t}.
    \end{aligned}
    $$

    $\exists~ \varphi(x) \in C^\infty(\mathbb{R})$, define $\varphi_n(x) = \begin{cases}
        -1, &x \leq 0\\
        \varphi(nx), & 0 < x < \frac{1}{n}\\
        1,& x \geq \frac{1}{n},
    \end{cases} $
    then $\varphi_n(x) \to \mathrm{sgn}(x)$, $\forall~ x \neq 0$.

    Let $$f_n(x) = \begin{cases}
        -x, & x \leq0\\
        \int_0^x \varphi_n(t) \diff t, & x >0.
    \end{cases}$$
    $$
    \implies f_n'(x) = \varphi_n(x.)
    $$
    Because $f(x) = \abs{ x } = \int_0^{\abs{ x }} 1 \diff t.$
    $$
    \implies \abs{ f(x) - f_n(x) } \leq \int_0^{\abs{ x }}\abs{ \varphi_n(t) - 1 }\diff  t \leq \int_0^{\frac{1}{n}}\abs{ \varphi_n(t)-1 }\diff  t \leq 2 \frac{1}{n}.
    $$

    Apply It\^o formula to $f_n$ to get that
    $$
    f_n(X_t) = f_n(X_0) + \int_0^t f_n'(X_s) \diff X_s + \frac{1}{2}\int_0^t f_n''(X_s) \diff [X]_s.
    $$
    Note
    $$
    \int_0^t f_n'(X_s) \diff X_s = \int_0^t f_n'(X_s) \diff M_s + \int_0^t f_n'(X_s) \diff A_s,
    $$
    and
    $$
    \int_0^t f_n'(X_s) \diff A_s = \int_0^t \varphi_n(X_s) \diff A_s.
    $$

    We have $\varphi_n(x) \uparrow \mathrm{sgn}(x)$, $\forall~x$. By Monotone convergence,
    $$
    \abs{ \int_0^t \varphi_n(X_s) - \mathrm{sgn}(X_s) \diff A_s } \leq \int_0^t \abs{ \varphi_n(X_s) - \mathrm{sgn}(X_s) } \abs{ \diff A_s } \to 0,
    $$
    which implies $\int_0^t \varphi_n(X_s)\diff A_s \to \int_0^t \mathrm{sgn}(X_s) \diff A_s$ uniform in $t$.

    Next, $T_k = \inf\{t \geq 0 \colon \abs{ M_t }\geq k\}$. Then $M^{T_k}$ is bounded. 
    
    We may assume $M$ is bounded. Then
    $$
    \norm{\int_0^\infty \left( \mathrm{sgn}(X_s) - f_n'(X_s) \right) \diff M_s}_2^2 = \mathbb{E}\int_0^\infty \left( \mathrm{sgn}(X_s) - f_n'(X_s) \right)^2 \diff [M]_s \to 0.
    $$

    So $$
    \mathbb{E}\left[ \sup_{t \leq T} \abs{ \int_0^t \left( \mathrm{sgn}(X_s) - f_n'(X_s) \right) \diff M_s}^2  \right] \xlongrightarrow{n \to \infty} 0.
    $$

    Hence 
    $$
    \sup_{t \geq 0} \int_0^t \left( \mathrm{sgn}(X_s) - f_n'(X_s) \right) \diff M_s \xlongrightarrow{a.s.} 0.
    $$
    $\implies \int_0^t f_n'(X_s) \diff X_s \xlongrightarrow{a.s.} \int_0^t \mathrm{sgn}(X_s) \diff  X_s$, $f_n(X_t) \xlongrightarrow{a.s.} \abs{ X_t }$, $f_n(X_0) \xlongrightarrow{a.s.} \abs{ X_0 }$.
    
    $\implies C_t^n \triangleq \int_0^t \frac{1}{2} f_n''(X_s) \diff [X_s]$, $C_t^n \xlongrightarrow{a.s.} l_t \triangleq \abs{ X_t } - \abs{ X_0 } - \int_0^t \mathrm{sgn}(X_s) \diff X_s.$

    $\diff C_s^n = \frac{1}{2}f_n''(X_s) \diff [X]_s \implies \forall~ \delta > 0$, if $\frac{1}{n} < \delta$, then
    $$
        \abs{ X_s } > \delta > \frac{1}{n} \implies f_n''(X_s) = 0.
    $$
    $$
    \implies\int_0^t \mathsf{1}_{\{\abs{ X_s } > \delta\}} \diff C_s^n = 0.
    $$

    Hence $$\int_0^t \mathsf{1}_{\{\abs{ X_s }> \delta\}} \diff  l_s \leq \lim_{n \to \infty} \int_0^t \mathsf{1}_{\{\abs{ X_s }> \frac{\delta}{2}\}} \diff C_s^n = 0.$$

    $$
    \int_0^t \varphi_\delta(X_s) \diff l_s = \lim_{n\to \infty}\int_0^t \varphi_\delta(X_s) \diff C_s^n \leq \liminf_{n \to \infty} \int_0^t \mathsf{1}_{\{\abs{ X_s } > \frac{\delta}{2}\}} \diff  C_s^n.
    $$

    We prove $\int_0^t \mathsf{1}_{\{\abs{ X_s }> \delta\}}\diff l_s = 0$, $\forall~ \delta >0$ $\implies \int_0^t \mathsf{1}_{\{\abs{ X_s } \neq 0\}} \diff  l_s =0.$
} 

Let $X,Y$ be continuous semimartingale. Then 
$$
\int_0^t Y_t \diff X_t = \text{It\^o Integral.}
$$

Stratonovich Integral
$$
\int_0^t Y_t \partial X_t \triangleq \int_0^t Y_t \diff X_t + \frac{1}{2}[X,Y]_{t}.
$$

Recall that 
$$
\begin{aligned}
X_t Y_t - X_0 Y_0 &= \int_0^t X_s \diff Y_s + \int_0^t Y_s \diff X_s + [X,Y]_t \\&= \int_0^t X_s \diff Y_s + \frac{1}{2}[X,Y]_t + \int_0^tY_s \diff X_s + \frac{1}{2}[X,Y]_t\\
&=\int_0^t X_s \partial Y_s + \int_0^t Y_s \partial X_s.
\end{aligned}
$$

\thmnn{
    $f(X_t) = f(X_0) + \int_0^t f'(X_s) \partial X_s$.
}

\pf{
    By It\^o formula,
    $$
    f(X_t) = f(X_0) + \int_0^t f'(X_s) \diff X_s + \frac{1}{2} \int_0^t f''(X_s) \diff [X_s].
    $$

    By definition, 
    $$
    \int_0^t f'(X_s) \partial X_s = \int_0^t f'(X_s) \diff X_s + \frac{1}{2}\left[ f'(X),X \right]_t.
    $$
    $$
    \diff f(X_t) = f'(X_t) \diff X_t + \frac{1}{2}f''(X_t) \diff [X]_t
    $$
    $$
    \implies \diff f'(X_t) = f''(X_t) \diff X_t + \frac{1}{2} f'''(X_t) \diff [X]_t.
    $$
    So
    $$
    f'(X_t) = f'(X_0) + \int_0^t f''(X_s) \diff X_s  + \frac{1}{2}\int_0^t f'''(X_s) \diff [X]_s.
    $$
    $$
    \implies \left[ f'(X),X \right]_t = \left[ \int_0^t f''(X_s) \diff X_s, X_t \right] = \int_0^t f''(X_s) \diff [X]_s.
    $$
    ($[H\cdot M, K \cdot N]_t = \int_0^t HK \diff [M,N]_s$).

    So 
    $$
    f(X_t) = f(X_0) + \int_0^t f'(X_s) \diff X_s + \frac{1}{2} \left[ f'(X), X \right]_t = f(x_0) + \int_0^t f'(X_s) \partial X_s.
    $$
}

\thm{Riemam-Sum Approximation}{
    Let $C$ be a locally bounded adapted $L$-process. Let $X$ be a continuous semimartingale. Then
    $$
    \sum_{i=1}^{\infty}c(t_{i-1}^n)(X_{t_i^n} - X_{t_{i-1}^n}) \to \int_0^t c(s) \diff X_s.
    $$
    with $t_i^n = t \wedge \frac{i}{2^n}$. The convergence is uniform in probability on compacts, i.e. $\forall~ K > 0$,
    $$
    \sup_{t \leq K}\abs{ I^n(t)-I(t) } \xlongrightarrow{P} 0.
    $$
}

\pf{
    Let $c^n(t) = c(t_{i-1}^n)$ for $t \in (t_{i-1}^n,t_i^n]$. Then $c^n(t) \to c(t)$ for all $t \geq 0.$

    Also, 
    $$
    \begin{aligned}
    I^n(t) &= \sum_{i=1}^\infty c(t_{i-1}^n)(X(t_{i}^n) - X(t_{i-1}^n))\\
    &= \int_0^t c^n(s) \diff X_s \\
    &= (c^n,X)_t = \int_0^t c^n(s)\diff M_s + \int_0^t c^n(s) \diff A_s.
    \end{aligned}$$

    $\forall~ t \leq k$, $\abs{ c^n(t) } \leq c_k$
    $$
    \implies\int_0^t c^n(s) \diff A_s \to \int_0^t c(s)\diff A_s.
    $$

    By localization, we may assume $M$ is bounded. If $c$ is bounded and $M \in \mathcal{M}_0^2$, then 
    $$
    \int_0^t c^n(s) \diff M_s \to \int_0^t c(s) \diff M_s.
    $$
    $$
    \mathbb{E}\left[ \sup_{t \leq k}\abs{ \int_0^t (c^n(s) - c(s)) \diff M_s }^2 \right] \leq c_k \mathbb{E}\int_0^k (c^n(s) - c(s))^2 \diff [M]_s \to 0.
    $$
}